\chapter{Projekt rozwiązania}

\section{Wymagania}
\subsection*{Wymagania funkcjonalne}
\begin{itemize}
    \item System musi umożliwiać równoległe uruchomienie wersji stabilnej SMF v1 oraz wersji shadow SMF v2.
    \item Ruch produkcyjny musi być obsługiwany przez SMF v1, a kopia żądań przekazywana do SMF v2.
    \item System musi wymuszać izolację stanu danych dla wersji shadow (oddzielna baza \texttt{udr-shadow}).
    \item Architektura musi umożliwiać ocenę jakości wersji v2 na podstawie metryk z Service Mesh.
\end{itemize}

\subsection*{Wymagania niefunkcjonalne}
\begin{itemize}
    \item Rozwiązanie musi być oparte o narzędzia \emph{cloud native} i opisywane jako IaC.
    \item Konfiguracja powinna być przenośna między środowiskami laboratoryjnymi i operatorskimi.
    \item Walidacja ma być bezinwazyjna dla ruchu użytkowników i odporna na awarie shadow.
    \item Proces oceny wdrożenia powinien wspierać automatyzację decyzji \emph{go/no-go}.
\end{itemize}

\section{Architektura logiczna}
Zaproponowana architektura obejmuje cztery warstwy:
\begin{enumerate}
    \item \textbf{Warstwa uruchomieniowa} -- klaster Kubernetes z namespace \texttt{5gc}~\cite{kubernetesDocs}.
    \item \textbf{Warstwa komunikacyjna} -- Istio z sidecarami Envoy oraz regułami routingu~\cite{istioDocs,envoyDocs}.
    \item \textbf{Warstwa funkcjonalna} -- wdrożenia SMF v1 i SMF v2.
    \item \textbf{Warstwa danych i walidacji} -- UDR produkcyjny, UDR shadow, Prometheus i Iter8.
\end{enumerate}

Model ruchu działa następująco: żądanie klienta trafia do usługi \texttt{smf}, następnie \texttt{VirtualService} przekazuje je do subsetu \texttt{v1}, a jego kopię do subsetu \texttt{v2}. Odpowiedź do klienta zwracana jest wyłącznie z v1.

\section{Projekt izolacji danych}
W pracy przyjęto zasadę twardej separacji stanu pomiędzy ścieżką produkcyjną i testową:
\begin{itemize}
    \item SMF v1 korzysta wyłącznie z \texttt{udr-prod},
    \item SMF v2 korzysta wyłącznie z \texttt{udr-shadow},
    \item poświadczenia dostępowe są rozdzielone na osobne sekrety Kubernetes.
\end{itemize}

Takie podejście chroni profile abonentów przed modyfikacją przez kod testowy i umożliwia bezpieczne uruchamianie eksperymentów na ruchu produkcyjnym.

\section{Projekt routingu Service Mesh}
Routing zrealizowano przez dwa obiekty Istio~\cite{istioDestinationRule,istioVirtualService}:
\begin{itemize}
    \item \texttt{DestinationRule} definiujący subsety \texttt{v1} oraz \texttt{v2} na podstawie etykiety \texttt{version},
    \item \texttt{VirtualService} realizujący:
    \begin{itemize}
        \item \texttt{route} 100\% do \texttt{v1},
        \item \texttt{mirror} do \texttt{v2},
        \item \texttt{mirrorPercentage} ustawione początkowo na 10\%.
    \end{itemize}
\end{itemize}

Taka konfiguracja pozwala zwiększać wolumen mirrorowanego ruchu etapowo, zależnie od wyników walidacji.

\section{Projekt obserwowalności i kryteriów SLO}
Walidacja jakości opiera się o metryki Envoy eksportowane do Prometheus~\cite{prometheusDocs}:
\begin{itemize}
    \item średnie opóźnienie żądań dla \texttt{smf-v2},
    \item odsetek błędów HTTP 4xx/5xx dla \texttt{smf-v2}.
\end{itemize}

Dla automatycznej oceny przyjęto progi:
\begin{itemize}
    \item \texttt{error\_rate <= 0},
    \item \texttt{latency\_avg\_ms < 100}.
\end{itemize}

Ocena jest realizowana przez eksperyment Iter8 (\texttt{custommetrics} + \texttt{assess}), co umożliwia wpięcie procesu w pipeline CI/CD~\cite{iter8Docs}.

\section{Artefakty projektowe}
Projekt został odwzorowany w artefaktach IaC:
\begin{itemize}
    \item Terraform: definicja bazowej warstwy platformowej klastra i namespace,
    \item YAML: workloady SMF, UDR, konfiguracja Istio oraz walidacja Iter8,
    \item YAML (opcjonalnie): deployment Diffy do porównania odpowiedzi.
\end{itemize}

Przyjęty podział plików umożliwia niezależne rozwijanie warstwy infrastruktury, routingu i walidacji jakości.

\section{Modelowanie C4 i materiał diagramowy}
Architektura została dodatkowo opisana zestawem diagramów C4 w składni Mermaid. Pełny kod źródłowy diagramów (C1, C2, C3, C4) umieszczono w dodatku \ref{chap:c4-mermaid}, natomiast instrukcję eksportu i osadzania wykresów w LaTeX przedstawiono w dodatku \ref{chap:mermaid-latex}.